\documentclass[a4paper,10pt,fleqn]{article}
\usepackage[margin=1in]{geometry}
\usepackage{amssymb}
\usepackage{adjustbox}
\usepackage{natbib}
\usepackage[colorlinks=true,linkcolor=blue,citecolor=blue,urlcolor=blue]{hyperref}
\usepackage{fuzz}
\usepackage{schemapk}
\usepackage{zed-maths}
\usepackage{zed-proof}
\newdimen\savedleftskip
\begin{document}

\section*{PURETEXT Block Examples}
\addcontentsline{toc}{section}{PURETEXT Block Examples}

\subsection*{Example 1 : Basic PURETEXT}
\addcontentsline{toc}{subsection}{Example 1 : Basic PURETEXT}

\bigskip

This is a verbatim text block. Unlike TEXT blocks, PURETEXT does not apply smart quote conversion or other typographical transformations.

\bigskip

\bigskip

The characters " and ' remain exactly as typed, without conversion to opening and closing quotes.

\bigskip

\subsection*{Example 2 : Preserving Special Characters}
\addcontentsline{toc}{subsection}{Example 2 : Preserving Special Characters}

\bigskip

In PURETEXT blocks, special characters like \& \% \$ \# \_ \{ \} are preserved literally without requiring LaTeX escaping.

\bigskip

\bigskip

Mathematical notation like x\textasciicircum{}2 or f(n) is also treated as literal text, not as mathematical expressions.

\bigskip

\subsection*{Example 3 : Code Snippets}
\addcontentsline{toc}{subsection}{Example 3 : Code Snippets}

\bigskip

When documenting code or command-line syntax, use PURETEXT to preserve exact formatting:

\bigskip

\bigskip

\$ txt2tex examples/hello\_world.txt

\bigskip

\bigskip

The above command compiles the example to PDF.

\bigskip

\subsection*{Example 4 : Literal Z Notation Syntax}
\addcontentsline{toc}{subsection}{Example 4 : Literal Z Notation Syntax}

\bigskip

To show Z notation syntax without rendering it as formal notation, use PURETEXT:

\bigskip

\bigskip

forall x : N | x > 0

\bigskip

\bigskip

This appears as literal text rather than being typeset as a quantified predicate.

\bigskip

\subsection*{Example 5 : Comparison TEXT vs PURETEXT}
\addcontentsline{toc}{subsection}{Example 5 : Comparison TEXT vs PURETEXT}

\noindent This block uses smart quotes: "Hello world" becomes properly formatted.

\bigskip

\bigskip

This block preserves quotes: "Hello world" stays exactly as typed.

\bigskip

\noindent Mathematical notation x\textasciicircum$\{\}$2 is rendered as a superscript.

\bigskip

\bigskip

Mathematical notation x\textasciicircum{}2 appears as literal text.

\bigskip

\subsection*{Example 6 : File Paths land URLs}
\addcontentsline{toc}{subsection}{Example 6 : File Paths land URLs}

\bigskip

File paths like /Users/name/Documents/file.txt are preserved exactly.

\bigskip

\bigskip

URLs like https://example.com/path?query=value\&other=123 remain unmodified.

\bigskip

\subsection*{Example 7 : Mixed Usage}
\addcontentsline{toc}{subsection}{Example 7 : Mixed Usage}

\noindent You can freely mix TEXT and PURETEXT blocks. Use TEXT for normal prose with smart typography.

\bigskip

\bigskip

Use PURETEXT when you need exact character-for-character reproduction.

\bigskip

\noindent This combination gives you flexibility: elegant typography when you want it, literal text when you need it.

\bigskip

\end{document}